\documentclass[UTF8]{ctexart}
\usepackage{amsmath,amssymb}
\usepackage[a4paper,margin=2.5cm]{geometry}
\title{高等代数 1022}
\date{}
\begin{document}
\maketitle

\section*{高等代数 1022}

\subsection*{4. 求和号}
\[
\sum_{i=1}^n a_i=a_1+\cdots+a_n.
\]

双重求和：
\[
\begin{array}{cccc@{\qquad}c}
a_{11} & a_{12} & \cdots & a_{1n}
  & \displaystyle\sum_{j=1}^{n} a_{1j} \\[2.5ex]
a_{21} & a_{22} & \cdots & a_{2n}
  & \displaystyle\sum_{j=1}^{n} a_{2j} \\[2.5ex]
\vdots & \vdots &        & \vdots
  & \vdots \\[2.5ex]
a_{m1} & a_{m2} & \cdots & a_{mn}
  & \displaystyle\sum_{j=1}^{n} a_{mj} \\[3.5ex]
\displaystyle\sum_{i=1}^{m} a_{i1}
  & \displaystyle\sum_{i=1}^{m} a_{i2}
  & \cdots
  & \displaystyle\sum_{i=1}^{m} a_{in}
  &
\end{array}
\]
如果想求整个数组的和，应表达为
\[
\sum_{i=1}^m\sum_{j=1}^n a_{ij}=\sum_{j=1}^n\sum_{i=1}^m a_{ij}.
\]

下三角求和：
\[
\sum_{i=1}^n\sum_{j=1}^i a_{ij}=\sum_{j=1}^n\sum_{i=j}^n a_{ij}.
\]

\subsection*{高次方程}

二次、三次、四次方程的一般形式：
\[
ax^2+bx+c=0,\qquad
ax^3+bx^2+cx+d=0,\qquad
ax^4+bx^3+cx^2+dx+e=0.
\]
五次及五次以上：Abel 定理、伽罗瓦（Galois）理论。

设 $K$ 是数域，$a_0,a_1,\cdots,a_n\in K$，$x$ 是字母，$a_0\neq0$，称
\[
a_0x^n+a_1x^{n-1}+\cdots+a_n=0
\]
是 $K$ 上的一元 $n$ 次方程。

如果 $x=\alpha\in\mathbb{C}$ 满足
\[
a_0\alpha^n+\cdots+a_n=0,
\]
则称 $\alpha$ 是方程的解或根。

代数基本定理：数域 $K$ 上的任何一个一元 $n$ 次方程在 $\mathbb{C}$ 上一定有解（有根）。

\end{document}
